% Part: normal-modal-logic
% Chapter: completeness
% Section: frame-completeness

\documentclass[../../../include/open-logic-section]{subfiles}

\begin{document}

\olfileid[pt]{nml}{com}{fra}

\olsection{Completude de frames}

O teorema da completude de \Log{K} pode ser estendido a outros
sistemas modais, uma vez que mostremos que o modelo canônico de uma
lógica dada tem a propriedade de frames correspondente.

\begin{thm}\ollabel{thm:completeframeprops}
  Se uma lógica modal normal $\Sigma$ contém uma das !!{formula}s do
  lado esquerdo de \olref{tab:correspondencetable}, então o modelo
  canônico de~$\Sigma$ tem a propriedade correspondente do lado
  direito.
\end{thm}

\begin{table}[htp]
  \centering
    \begin{tabular}{| l || l |}
      \hline
      \emph{Se $\Sigma$ contém \dots} & \emph{\dots o modelo canônico
        de $\Sigma$ é:} \\
      \hline \hline
      \Ax{D}:\;\; $\Box!A \lif \Diamond !A$ & \quad serial; \\
      \hline
      \Ax{T}:\;\; $\Box!A \lif !A$ &\quad  reflexivo;\\
      \hline
      \Ax{B}: \;\;$!A \lif \Box\Diamond!A$ &\quad  simétrico; \\
      \hline
      \Ax{4}: \;\; $\Box!A \lif \Box\Box!A$ & \quad transitivo; \\
      \hline 
      \Ax{5}: \;\; $\Diamond !A \lif \Box\Diamond!A$& \quad euclidiano.\\
      \hline
    \end{tabular}
    \caption{Fatos básicos de correspondência.}
    \ollabel{tab:correspondencetable}
\end{table}

\begin{proof}
  Tratamos de cada uma dessas propriedades sucessivamente.

  Suponha que $\Sigma$ contenha \Ax{D} e seja
  $\Delta \in W^\Sigma$; devemos mostrar que existe um $\Delta'$ tal
  que $R^\Sigma \Delta\Delta'$. Basta mostrar que
  $\Box^{-1}\Delta$ é $\Sigma$-consistente, pois então, pelo lema de
  Lindenbaum, existe um conjunto $\Sigma$-consistente completo
  $\Delta' \supseteq \Box^{-1}\Delta$ e, pela definição de $R^\Sigma$
  \iftag{prvBox}{}{e por \olref[mod]{lem:box-iff-diamond} }, temos
  $R^\Sigma \Delta\Delta'$. Suponha, portanto, por contradição, que
  $\Box^{-1}\Delta$ \emph{não} seja $\Sigma$-consistente, isto é,
  $\Box^{-1}\Delta \Proves[\Sigma] \lfalse$. Por \olref[mod]{lem:box2},
  $\Delta \Proves[\Sigma] \Box\lfalse$ e, como $\Sigma$ contém
  \Ax{D}, também $\Delta \Proves[\Sigma] \Diamond\lfalse$. Mas $\Sigma$
  é normal, logo $\Sigma \Proves \lnot\Diamond \lfalse$
  (\olref[prf][nor]{prop:notDiamondBot}); daí também $\Delta
  \Proves[\Sigma] \lnot\Diamond \lfalse$, contrariando a consistência
  de~$\Delta$.

  Suponha agora que $\Sigma$ contenha \Ax{T} e seja $\Delta \in
  W^\Sigma$. Queremos mostrar $R^\Sigma \Delta\Delta$, isto
  é,\iftag{prvBox}{}{ por \olref[mod]{lem:box-iff-diamond},}
  $\Box^{-1}\Delta \subseteq \Delta$. Mas, se $\Box!A \in \Delta$,
  então, por \Ax{T}, também $!A \in \Delta$, como desejado.

  Suponha agora que $\Sigma$ contenha \Ax{B}, e suponha
  $R^\Sigma \Delta\Delta'$ para $\Delta$, $\Delta' \in W^\Sigma$.
  Devemos mostrar que $R^\Sigma \Delta'\Delta$, isto é,\iftag{prvBox}{
    $\Box^{-1}\Delta' \subseteq \Delta$. Por
    \olref[mod]{lem:box-iff-diamond}, isso é equivalente a}{}
  $\Diamond\Delta \subseteq \Delta'$. Suponha, portanto,
  $!A \in \Delta$. Por \Ax{B}, também $\Box\Diamond!A \in \Delta$.
  Pela hipótese de que $R^\Sigma \Delta\Delta'$\iftag{prvBox}{}{ e por
    \olref[mod]{lem:box-iff-diamond}}, temos $\Box^{-1}\Delta
  \subseteq \Delta'$ e, portanto, $\Diamond!A \in \Delta'$, como exigido.

  Suponha agora que $\Sigma$ contenha \Ax{4}, e suponha
  $R^\Sigma \Delta_1\Delta_2$ e $R^\Sigma \Delta_2\Delta_3$. Devemos
  mostrar $R^\Sigma \Delta_1\Delta_3$. Da hipótese\iftag{prvBox}{}{ e de
    \olref[mod]{lem:box-iff-diamond}}, temos tanto
  $\Box^{-1}\Delta_1 \subseteq \Delta_2$ quanto
  $\Box^{-1}\Delta_2 \subseteq \Delta_3$. Para mostrar
  $R^\Sigma \Delta_1\Delta_3$, basta mostrar
  $\Box^{-1}\Delta_1 \subseteq \Delta_3$. Seja, portanto,
  $!B \in \Box^{-1}\Delta_1$, isto é, $\Box!B \in \Delta_1$. Por
  \Ax{4}, também $\Box\Box!B \in \Delta_1$ e, pela hipótese, obtemos,
  primeiro, $\Box!B \in \Delta_2$ e, segundo, $!B \in \Delta_3$, como
  desejado.

  Suponha agora que $\Sigma$ contenha \Ax{5}, e suponha
  $R^\Sigma \Delta_1\Delta_2$ e $R^\Sigma \Delta_1\Delta_3$. Devemos
  mostrar $R^\Sigma \Delta_2\Delta_3$. A primeira hipótese fornece
  $\Box^{-1}\Delta_1 \subseteq \Delta_2$\iftag{prvBox}{}{ por
    \olref[mod]{lem:box-iff-diamond}}, e a segunda hipótese é
  equivalente a $\Diamond\Delta_3 \subseteq \Delta_1$\iftag{prvBox}{,
    por \olref[mod]{lem:box-iff-diamond}}{}. Para mostrar $R^\Sigma
  \Delta_2\Delta_3$\iftag{prvBox}{, por
    \olref[mod]{lem:box-iff-diamond}, basta}{, devemos}
  mostrar $\Diamond\Delta_3 \subseteq \Delta_2$. Seja
  $\Diamond!A \in \Diamond\Delta_3$, isto é, $!A \in \Delta_3$. Pela
  segunda hipótese, $\Diamond!A \in \Delta_1$ e, por \Ax{5}, também
  $\Box\Diamond!A \in \Delta_1$. Mas agora a primeira hipótese fornece
  $\Diamond!A \in \Delta_2$, como desejado.
  \end{proof}

Como corolário, obtemos resultados de completude para vários sistemas.
Por exemplo, sabemos que $\Log{S5} = \Log{KT5} = \Log{KTB4}$ é
completa com respeito à classe de todos os modelos reflexivos e
euclidianos, que é a mesma classe de todos os modelos reflexivos,
simétricos e transitivos.

\begin{thm}\ollabel{thm:generaldet}
  Sejam $\mClass{C}_\Ax{D}$, $\mClass{C}_\Ax{T}$,
  $\mClass{C}_\Ax{B}$, $\mClass{C}_\Ax{4}$ e
  $\mClass{C}_\Ax{5}$ as classes de todos os modelos seriais,
  reflexivos, simétricos, transitivos e euclidianos, respectivamente.
  Então, para quaisquer esquemas $!A_1$, \dots, $!A_n$ entre \Ax{D},
  \Ax{T}, \Ax{B}, \Ax{4} e \Ax{5}, o sistema
  $\Log{K}!A_1 \dots !A_n$ é determinado pela classe de modelos
  $\mClass{C} = \mClass{C}_{!A_1} \cap \dots
  \cap \mClass{C}_{!A_n}$. 
\end{thm}

\begin{prop}
  Seja $\Sigma$ uma lógica modal normal; então:
  \begin{enumerate}
  \item\ollabel{prop:anotherfive-a}%
    Se $\Sigma$ contém o esquema $\Diamond!A \lif \Box
    !A$, então o modelo canônico de $\Sigma$ é parcialmente funcional.
  \item Se $\Sigma$ contém o esquema $\Diamond!A \liff \Box
    !A$, então o modelo canônico de $\Sigma$ é funcional.
  \item Se $\Sigma$ contém o esquema $\Box\Box!A \lif \Box
    !A$, então o modelo canônico de $\Sigma$ é fracamente denso.
  \end{enumerate}
(veja \olref[frd][acc]{tab:anotherfive} para as definições dessas
  propriedades de frames).
\end{prop}

\begin{proof}
  \begin{enumerate}
  \item Suponha que $\Sigma$ contenha o esquema $\Diamond !A \lif
    \Box !A$. Para mostrar que $R^\Sigma$ é parcialmente funcional,
    devemos provar que, para quaisquer $\Delta_1$, $\Delta_2$,
    $\Delta_3 \in W^\Sigma$, se $R^\Sigma \Delta_1\Delta_2$ e
    $R^\Sigma \Delta_1\Delta_3$, então $\Delta_2=\Delta_3$. Como
    $R^\Sigma \Delta_1\Delta_2$, temos
    $\Box^{-1}\Delta_1 \subseteq \Delta_2$ e, como
    $R^\Sigma \Delta_1\Delta_3$, também
    $\Box^{-1}\Delta_1 \subseteq \Delta_3$\iftag{prvBox}{}{, ambos por
      \olref[mod]{lem:box-iff-diamond}}. A identidade
    $\Delta_2=\Delta_3$ decorrerá se pudermos estabelecer as duas
    inclusões $\Delta_2 \subseteq \Delta_3$ e
    $\Delta_3 \subseteq \Delta_2$. Para a primeira inclusão, seja
    $!A \in \Delta_2$; então $\Diamond!A \in \Delta_1$ e, pelo esquema
    e pelo fechamento dedutivo de $\Delta_1$, também
    $\Box!A \in \Delta_1$; daí, pela hipótese de que
    $R^\Sigma \Delta_1\Delta_3$, $!A \in \Delta_3$. A segunda inclusão
    é semelhante.

  \item Isso decorre imediatamente da parte~\olref{prop:anotherfive-a}
    e da demonstração da serialidade em \olref{thm:completeframeprops}.

  \item Suponha que $\Sigma$ contenha o esquema
    $\Box\Box!A \lif \Box!A$ e, para mostrar que $R^\Sigma$ é fracamente
    densa, seja $R^\Sigma \Delta_1\Delta_2$. Devemos mostrar que existe
    um conjunto $\Sigma$-consistente completo $\Delta_3$ tal que
    $R^\Sigma \Delta_1\Delta_3$ e $R^\Sigma \Delta_3\Delta_2$. Seja:
    \[
    \Gamma = \Box^{-1}\Delta_1 \cup \Diamond\Delta_2.
    \]
    Basta mostrar que $\Gamma$ é $\Sigma$-consistente, pois então,
    pelo lema de Lindenbaum, ele pode ser estendido a um conjunto
    $\Sigma$-consistente completo~$\Delta_3$ tal que
    $\Box^{-1}\Delta_1 \subseteq \Delta_3$ e
    $\Diamond\Delta_2 \subseteq \Delta_3$, isto é,
    $R^\Sigma \Delta_1\Delta_3$\iftag{prvBox}{}{ (por
      \olref[mod]{lem:box-iff-diamond})} e $R^\Sigma
    \Delta_3\Delta_2$\iftag{prvBox}{ (por
      \olref[mod]{lem:box-iff-diamond})}{}.

    Suponha, por contradição, que $\Gamma$ não seja consistente. Então
    existem !!{formula}s $\Box!A_1$, \dots, $\Box!A_n \in \Delta_1$ e
    $!B_1$, \dots,~$!B_m \in \Delta_2$ tais que \[!A_1, \dots,
    !A_n, \Diamond!B_1, \dots, \Diamond!B_m \Proves[\Sigma]
    \lfalse.\] Como $\Diamond (!B_1 \land \dots \land !B_m) \to
    (\Diamond!B_1 \land \dots \land \Diamond!B_m)$ é !!{derivable} em
    toda lógica modal normal, argumentamos como segue, contradizendo a
    consistência de~$\Delta_2$:
    \begin{align*}
      !A_1, \dots, !A_n, & \Diamond!B_1, \dots,
      \Diamond!B_m \Proves[\Sigma] \lfalse \\
      !A_1, \dots,!A_n & \Proves[\Sigma] (\Diamond!B_1 \land
      \dots \land \Diamond!B_m) \lif \lfalse\\
      & \qquad\text{pelo teorema da dedução}\\
      & \qquad\text{\olref[prf][prp]{prop:derivabilityfacts}\olref[prf][prp]{prop:derivabilityfacts-deduction}, e \Taut} \\
      !A_1, \dots,!A_n & \Proves[\Sigma]
      \Diamond(!B_1 \land
      \dots \land !B_m) \lif \lfalse\\
      & \qquad \text{pois $\Sigma$ é normal} \\
      !A_1, \dots,!A_n & \Proves[\Sigma]
      \lnot\Diamond (!B_1 \land \dots \land !B_m)\\
      & \qquad\text{por \PL} \\
      !A_1, \dots,!A_n & \Proves[\Sigma]
      \Box\lnot (!B_1 \land \dots \land !B_m)\\
      & \qquad\text{$\Box\lnot$ por $\lnot\Diamond$} \\
      \Box!A_1, \dots,\Box!A_n
      & \Proves[\Sigma] \Box\Box \lnot (!B_1 \land \dots \land !B_m)\\
      & \qquad\text{por \olref[mod]{lem:box1}} \\
      \Box!A_1, \dots,\Box!A_n
      & \Proves[\Sigma]
      \Box\lnot (!B_1 \land \dots \land !B_m)\\
      &\qquad\text{pelo esquema $\Box\Box!A \lif \Box!A$} \\
      \Delta_1 & \Proves[\Sigma]
      \Box\lnot (!B_1 \land \dots \land !B_m)\\
      &\qquad\text{por monotonicidade, \olref[prf][prp]{prop:derivabilityfacts}\olref[prf][prp]{prop:derivabilityfacts-monotonicity}} \\
      & \Box\lnot (!B_1 \land \dots \land !B_m) \in \Delta_1\\
      &\qquad\text{pelo fechamento dedutivo}; \\
      & \lnot (!B_1 \land \dots \land !B_m) \in \Delta_2\\
      &\qquad \text{pois }
      R^\Sigma \Delta_1\Delta_2. 
    \end{align*}
  \end{enumerate}
\end{proof}

Com base nesses exemplos, poderíamos pensar que todo sistema~$\Sigma$
de lógica modal é \emph{completo}, no sentido de demonstrar toda
fórmula válida em todo frame no qual todo teorema de $\Sigma$ é
válido. Infelizmente, há muitos sistemas que não são completos nesse
sentido.

\end{document}
